\section*{Supplemental Methods}

\subsection*{Assembly and annotation of \textit{Balanus crenatus}}\label{sec:sup_methods_balcre}

\subsubsection*{High-molecular weight DNA and PacBio sequencing}

For genome sequencing using PacBio Hi-Fi reads, we sampled several \textit{Balanus
crenatus} individuals from the docks in the Oregon Institute of Marine Biology, in
Charleston, Oregon, USA (43.346 N, 124.325 W) in November 2023 during low tide.
High-molecular weight (HMW) DNA was extracted from the dissected testes using the 
PacBio Nanobind tissue kit (PacBio).
DNA was assessed for quality by fluorometric quantification of concentration (Qubit;
Invitrogen) and fragment length distribution (Fragment Analyzer; Advanced Analytical).
The extracted DNA of a single individual was sequenced on two SMRTcells of a PacBio
Sequel II machine at the University of Oregon's Genomics and Cell Characterization
Core Facility (GC3F).
These two sequencing runs yielded a total of 3.06 million reads with a read N50 of
14.8\,kbp.

\subsubsection*{Estimating genome size and heterozygosity}

We used k-mers to estimate the size and heterozygosity of the \textit{B.
crenatus} genome.
First, we counted k-mers present in the PacBio HiFi reads using
\textsc{jellyfish} version \texttt{2.2.10} \citep{marcais_fast_2011}, counting
21-mers (\texttt{--mer-len 21}) present in both strands (\texttt{--canonical}).
These counts were then used to generate an empirical distribution of 21-mers
using the \textsc{jellyfish} \texttt{histo} command.
The final model of the diploid, genome-wide k-mer distribution was then performed
using \textsc{GenomeScope2} version \texttt{2.0}
\citep{ranallo-benavidez_genomescope_2020}.
We estimated the \textit{B. crenatus} genome size to be 755{,}696{,}144\,bp in length,
with a heterozygosity of 2.25\%.

\subsubsection*{Contig-level genome assembly}

We generated a contig-level genome assembly using \textsc{hifiasm} version
\texttt{0.19.8-r603} \citep{cheng_haplotype-resolved_2021,cheng_haplotype-resolved_2022},
using strict parameters for the identification and purging of haplotig sequences.
This assembly was generated using an estimated haploid genome size of 800\,Mbp
(\texttt{--hg-size 800m}), dropping k-mers observed over ten times the average coverage
(\texttt{-D 10.0}), and allowing scaffolding (\texttt{--dual-scaf}).
For handling haplotigs, we set the expected coverage at homozygous sites to
40x (\texttt{--hom-cov 40}), the upper bound of coverage to purging duplicates to 40x
(\texttt{--purge-max 40}), and the similarity threshold between duplicate haplotigs
to 25\% (\texttt{-s 0.25}).
We assessed the contiguity of this assembly using \textsc{quast} version \texttt{5.2.0}
\citep{gurevich_quast_2013}, showing an assembly composed of 2{,}932 contigs, with a total
length of 1.20\,Gbp, and a contig N50 of 746\,kbp.
Additionally, \textsc{compleasm} version \texttt{0.12-r237} \citep{huang_compleasm_2023}
revealed this assembly 94.96\% complete (with 19.64\% duplicates) against the
\texttt{arthropoda\_odb10} reference ortholog dataset from \textsc{BUSCO}
\citep{simao_busco_2015,manni_busco_2021,tegenfeldt_orthodb_2025}.

\subsubsection*{Identifying contamination}

% TODO: @Scott, please verify this.
We screened the \textit{B. crenatus} contig-level assembly for contamination using
\textsc{MMseqs2} release \texttt{15\-6f452} \citep{steinegger_mmseqs2_2017}.
After indexing the assembly, taxonomic assignment was done for each contig.
As described in the \textsc{MMseqs2} documentation, protein fragments were extracted
across the six possible frames, which were then mapped against the NCBI \texttt{nr} and
\texttt{SwissProt} databases.
We retained only sequences assigned to the order Thecostraca (NCBI taxid
\texttt{116172}).
This process removed 236 fragments spanning $\approx$ 19.7\,Mbp, producing an assembly
with 2{,}696 contigs, a total size of 1.18\,Gbp, and a contig N50 of 758.4\,kbp.
The resulting gene-completeness against the \texttt{arthropoda\_odb10} reference
set was 94.57\% with 19.45\% duplicates.

\subsubsection*{Purging haplotigs}

After screening the assembly for contamination, we identified and purged haplotig
sequences using \textsc{purge\_dups} version \texttt{1.2.5}
\citep{guan_identifying_2020}.
We first mapped the PacBio HiFi reads to the contig-level assembly with
\textsc{minimap2} \texttt{-x map-hifi}.
From these alignments, we calculated the read-depth histogram using \texttt{pbcstat}
and determined the base-level coverage cutoffs using \texttt{calcuts}.
After splitting the assembly at gaps with \texttt{split\_fa}, we performed an
assembly self-alignment using \textsc{minimap2} \texttt{-x asm20}.
Haplotigs were then marked according to the empirical coverage cutoffs using the
\texttt{purge\_dups} command.
Lastly, the assembly FASTA was processed to remove sequences tagged as haplotigs using
\texttt{get\_seqs}, only removing sequences at the ends of contigs (\texttt{-e}) and
allowing the splitting of contigs (\texttt{-s}).
The purging reduced the number of contigs in the assembly to 1{,}548, resulting in
a total assembly size of 907.4\,Mbp, a contig N50 of 925.5\,kbp, and a gene-completeness
of 93.39\% with 3.16\% duplicates (against the \texttt{arthropoda\_odb10} reference
set).

\subsubsection*{Polishing and correcting}

The purged contig-level assembly was then polished and self-corrected using
\textsc{Inspector} software version \texttt{1.0.1} \citep{chen_accurate_2021}.
First, we evaluated the assembly using \texttt{inspector.py}, mapping the PacBio HiFi
reads against the assembly.
The assembly was then corrected using \texttt{inspector-correct.py}.
After polishing, the assembly contained 1{,}548 contigs, a total length of 907.3\,Mbp,
and a contig N50 of 925.6\,kbp.
Gene-completeness against the \texttt{arthropoda\_odb10} reference set was 93.39\%
with 3.26\% duplicates.

\subsubsection*{Reference-based scaffolding}

We used a \textsc{cactus} whole-genome alignment \citep{armstrong_progressive_2020} to
scaffold the assembly.
Using \textsc{cactus} version \texttt{2.9.7} we aligned the chromosome-level
\textit{B. glandula} assembly against the \textit{B. crenatus} contigs,
exporting the output alignments in MAF format.
After alignment, we used the \textsc{ragout} version \texttt{2.3} software
\citep{kolmogorov_ragout_2018} to perform reference-based scaffolding, disabling
the breaking of input sequences (\texttt{--solid-scaffolds}) and enabling the
repeat resolution algorithm (\texttt{--repeats}).
After running \textsc{ragout}, we only kept the \textit{B. crenatus} sequences that
were successfully aligned to \textit{B. glandula}.
This scaffolded assembly was composed of 220 sequence fragments, with a total length
of 851\,Mbp, and a scaffold N50 of 5.58\,Mbp.
Additionally, the assembly was 89.64\% gene-complete (with 2.37\% duplicates) when
compared to the \texttt{arthropoda\_odb10} reference set.

\subsubsection*{Genome annotation}

Repetitive elements were annotated in the scaffolded \textit{B. crenatus} assembly
using the \textsc{EarlGrey} software version \texttt{4.1.0} \citep{baril_earl_2024}.
When running \textsc{EarlGrey}, we provided the Crustacea \textsc{Dfam 3.8} partition
as an initial consensus library (\texttt{-l}), specified ``arthropoda'' as the
search term for \textsc{RepeatMasker} version \texttt{4.1.2}
\citep{smit_repeatmasker_2013}, clustered the final TE library to reduce redundancy
(\texttt{-c yes}), and generated a soft-masked FASTA as output (\texttt{-d yes}).
Similar to \textit{B. glandula}, the repeat annotation for \textit{B. crenatus}
revealed a highly repetitive genome, with 494.0\,Mbp (58.03\%) of the assembly
masked.
Moreover, the largest proportion of identified repeats remained unclassified,
289.4\,Mbp or 40.0\% of the assembly, further highlighting the lack of
barnacle-representative sequences in public repeat databases.
\vspace{0.2cm}

Following the annotation of repeats, we annotated protein-coding sequences in the
\textit{B. crenatus} assembly by lifting over gene-models from the curated \textit{B.
glandula} annotation using \textsc{LiftOn} version \texttt{1.0.5}
\citep{chao_lifton_2024}.
\textsc{LiftOn} lifts over annotations between genomes by integrating both
protein-to-genome alignments with \textsc{miniprot} \citep{li_miniprot_2023} and
nucleotide-to-nucleotide alignments using \textsc{Liftoff} \citep{shumate_liftoff_2021}.
Given the high diversity observed within these genomes and the high divergence
between them, we relaxed several of the \textsc{LiftOn} parameters to maximize the
DNA and protein sequence identity scores, as well as maintaining a high \textsc{BUSCO}
completeness.
We included 10\% of the flanking sequence of each gene as part of the alignment
(\texttt{-f flank 0.1}), we allowed the software to identify gene copies with sequence
similarity of at least 50\% (\texttt{-copies -sc 0.5}), used a 25\% coverage
threshold for the mapping of parent (\texttt{-a 0.25}) and child sequences
(\texttt{-s 0.25}), and enabled a post-liftover refinement of the transferred
gene models (\texttt{-polish}).
We also used the ``general'' splice model for \textsc{miniprot} (\texttt{-j 1}).
This liftover successfully annotated 15{,}841 genes, including 15{,}904 protein-coding
transcripts, in the \textit{B. crenatus} assembly.
While this is $\approx$ 5{,}000 fewer protein-coding genes than those annotated for
\textit{B. glandula}, this annotation was still 82.63\% gene-complete, with 9.67\%
duplicates, when compared to the \texttt{arthropoda\_odb10} reference ortholog set
using \textsc{compleasm} \texttt{protein}.

\subsection*{Processing annotations of published barnacle assemblies}\label{sec:sup_methods_barnacle_rna}

\subsubsection*{\textit{Pollicipes pollicipes}}

We used the published RefSeq annotation for the gooseneck barnacle \textit{P.
pollicipes} (NCBI accession: \texttt{GCF\_011947565.3};
\cite{bernot_chromosome-level_2022}; Table~\ref{tab:sup_ncbi_genomes}).
This annotation was processed with \textsc{agat} version \texttt{1.4.3}
\citep{dainat_another_2022} to filter out non-coding annotations, extract the longest 
transcript as representative for the gene, and standardize gene and sequence IDs for
downstream analysis.
This RefSeq annotation contains 20{,}420 protein-coding genes and is 88.72\%
gene-complete, with 35.93\% duplicates, when validated against the
\texttt{arthropoda\_odb10} reference ortholog dataset using \textsc{compleasm}
\texttt{protein}.

\subsubsection*{\textit{Amphibalanus amphitrite}}

We combined the genome and the annotation across two different assemblies for the
striped barnacle, \textit{A. amphitrite}.
First, we used the chromosome-level genome assembly for the species generated by
\cite{han_new_2024}.
While the assembled sequences for this genome are available (NCBI accession:
\texttt{GCA\_037642225.1}; Table~\ref{tab:sup_ncbi_genomes}), no annotation has
been published for this assembly.
Therefore, we then took the published RefSeq annotation (NCBI accession:
\texttt{GCF\_019059575.1}; Table~\ref{tab:sup_ncbi_genomes}), based on the
scaffold-level genome assembly by \cite{kim_draft_2019}.
We annotated the chromosome-level assembly by lifting over the RefSeq annotation using
the \textsc{LiftOn} software version \texttt{1.0.5} \citep{chao_lifton_2024},
allowing the software to identify gene copies with a sequence similarity of 80\%
(\texttt{-copies -sc 0.8}), and using the ``general'' splice model in
\textsc{miniprot} (\texttt{-j 1}).
This liftover process successfully annotated 23{,}035 protein-coding genes.
The annotation was then processed with \textsc{agat} version \texttt{1.4.3} to
extract the longest transcript as representative for the gene and to standardize gene
and sequence IDs for downstream analysis.
This processed annotation was then validated using \textsc{compleasm} \texttt{protein}
using the \texttt{arthropoda\_odb10} reference ortholog set, showing it was 88.08\%
complete with 12.24\% duplicates.

\subsubsection*{\textit{Amphibalanus improvisus}}

We annotated the chromosome-level assembly of the Bay barnacle, \textit{A. improvisus}
(NCBI accession: \texttt{GCA\_964274985.1};~\cite{bishop_genome_2025};
Table~\ref{tab:sup_ncbi_genomes}) using publicly available transcriptomic data.
We downloaded two public RNAseq datasets for this species, from NCBI BioProjects
\texttt{PRJNA528777} and \texttt{PRJNA528169} (Table~\ref{tab:sup_sra_data};
\cite{abramova_sensory_2019}).
The raw reads were then processed using \textsc{fastp} version \texttt{0.23.4}
\citep{chen_ultrafast_2023,chen_fastp_2018}, and aligned to the \textit{A. improvisus}
assembly using \textsc{hisat2} version \texttt{2.2.1} \citep{kim_graph-based_2019}.
We then used \textsc{BRAKER} version \texttt{3.0.8} \citep{gabriel_braker3_2024} to
annotate the genome using the evidence from both the RNAseq alignments and the
Arthropod representative protein sequences from \textsc{OrthoDBv11}
\citep{kriventseva_orthodb_2019,zdobnov_orthodb_2021}.
In \textsc{BRAKER}, we enforced the recovery of orthologs from the Arthropod
\textsc{BUSCO} dataset in the annotation (\texttt{--busco\_lineage=arthropoda\_odb10}).
After running \textsc{BRAKER}, we processed the resulting annotation using \textsc{agat}
to extract the longest transcript as representative for the gene and to standardize gene
and sequence IDs for downstream analysis.
This process annotated 16{,}098 protein-coding genes, and was 96.05\% gene-complete,
with 16.36\% duplicates, after validation with \textsc{compleasm} against the
\texttt{arthropoda\_odb10} reference ortholog dataset.

\subsubsection*{\textit{Capitulum mitella}}

We annotated the publicly-available chromosome-level genome assembly of the pedunculate
barnacle, \textit{C. mitella} (NCBI accession: \texttt{GCA\_030062745.1};
\cite{chen_capmit_2021}; Table~\ref{tab:sup_ncbi_genomes}) using RNAseq short read data
publicly available on the NCBI SRA database (Table~\ref{tab:sup_sra_data}).
We followed the same methodology used for \textit{A. improvisus}: processing the
short reads with \textsc{fastp}, aligning to the genome with \textsc{hisat2}, annotating
with \textsc{BRAKER}, and processing the resulting annotation with \textsc{agat}.
This approach resulted in the annotation of 10{,}504 protein-coding genes, which
produced an annotation that was 97.53\% gene-complete, with 14.02\% duplicate sequences
when compared to the \texttt{arthropoda\_odb10} reference ortholog dataset.
However, given the comparatively smaller number of annotated sequences, and the
resulting reduction in the number of orthologous sequences recovered when incorporating
this species, we excluded \textit{C. mitella} from most downstream analyses.

\subsection*{Codon usage biases in \textit{Drosophila melanogaster}}

As a comparison against barnacles, we calculated codon usage biases in the fruit
fly \textit{D. melanogaster}.
To do this, we first downloaded the reference annotation from NCBI RefSeq
(Release 6 plus ISO1 MT, NCBI accession: \texttt{GCF\_000001215.4};
Table~\ref{tab:sup_ncbi_genomes}).
We processed the annotation using \textsc{agat} version \texttt{1.4.3}
\citep{dainat_another_2022} to select the longest isoform per gene
and clean the IDs of the extracted coding sequences.
We calculated codon usage biases using the \textsc{cubar} R package version
\texttt{1.2.0} \citep{liu_cubar_2026}.
We filtered coding sequences using the \texttt{check\_cds} function, ensuring the
presence of canonical start and stop codons, retaining sequences for 13{,}521 genes.
Then we calculated the distribution of observed codons using the
\texttt{count\_codons} function.
Lastly, we calculated the effective number of codons (ENC) observed in the
empirical codon distribution using the \texttt{get\_enc} function.
This statistic determines the bias in codon usage in the form of the deviation from
equal codon usage across synonymous codons \citep{wright_effective_1990}, using the 
adjustments by~\cite{sun_improved_2013} to account for usage differences across
different codon subfamilies.

\subsection*{Nucleotide diversity in \textit{Drosophila melanogaster}}

We calculated nucleotide diversity ($\pi$) in sub-Saharan populations of \textit{D.
melanogaster} to establish a baseline for comparison against our central Oregon barnacle
population.
To achieve this, we first obtained the sequencing data for phase 2 of the Drosophila
Population Genomics Project (DPGP2;~\citep{pool_population_2012,lack_dpgp_2016}).
We processed the aligned FASTA files for all sub-Saharan samples (the `RG' population)
to generate a VCF using the \textsc{snp-sites} version \texttt{2.5.1} software
\citep{page_snp-sites_2016}.
We exported both variant and invariant sites (\texttt{-b}) for downstream compatibility.
The resulting gVCFs were then processed using a custom AWK command and \textsc{BCFtools}
version \texttt{1.21} \citep{danecek_bcftools_2021} to recode deletions, remove missing
alleles, and remove sites with over 25\% missing genotypes.
Nucleotide diversity along 10\,kbp windows was then calculated from the filtered VCF using
\textsc{pixy} version \texttt{2.0.0.beta12} \citep{korunes_pixy_2021,bailey_pixy2_2025},
allowing multiallelic sites to be used in the calculation
(\texttt{--include\_multiallelic\_snps}).

